\documentclass[11pt,a4paper]{article}
\usepackage[utf8]{inputenc}
\usepackage[T1]{fontenc}
\usepackage{amsmath, amssymb, amsthm}
\usepackage{geometry}
\geometry{margin=2.5cm}
\usepackage{lmodern}
\usepackage[final,expansion=false]{microtype}
\usepackage{booktabs}
\usepackage{array}
\usepackage{enumitem}
\usepackage{xcolor}
\usepackage{titlesec}
\PassOptionsToPackage{hyphens}{url}
\usepackage{hyperref}
\usepackage{xurl}
\hypersetup{
  colorlinks=true,
  linkcolor=blue!50!black,
  urlcolor=blue!50!black,
  citecolor=blue!50!black,
  pdftitle={Universality for orthoscheme reflection groups: the right-angled Coxeter envelope and the collision depth},
  pdfauthor={Vico Bonfioli},
}
\usepackage{parskip}
\setlength{\parindent}{0pt}
\setlength{\parskip}{0.55em}

\titleformat{\section}{\normalfont\Large\bfseries\color{blue!40!black}}{\thesection.}{0.6em}{}

\newcommand{\Z}{\mathbb{Z}}
\newcommand{\R}{\mathbb{R}}
\newcommand{\Q}{\mathbb{Q}}
\newcommand{\Isom}{\mathrm{Isom}}
\newcommand{\cd}{\mathrm{cd}}

\newtheorem{theorem}{Theorem}
\newtheorem{lemma}[theorem]{Lemma}
\newtheorem{corollary}[theorem]{Corollary}
\newtheorem{proposition}[theorem]{Proposition}
\newtheorem{conjecture}[theorem]{Conjecture}
\newtheorem{problem}[theorem]{Problem}
\theoremstyle{definition}
\newtheorem{definition}[theorem]{Definition}
\newtheorem{example}[theorem]{Example}
\theoremstyle{remark}
\newtheorem{remark}[theorem]{Remark}
\theoremstyle{plain}

\title{\textbf{Universality for orthoscheme reflection groups:\\
the right-angled Coxeter envelope and the collision depth}}
\author{Vico Bonfioli\thanks{Independent researcher, Italy.
\texttt{vicobonfioli@gmail.com}.\,
\url{https://elvec1o.github.io/home/posts/}.}}
\date{\today}

\begin{document}
\sloppy
\maketitle

\begin{abstract}
The $n+1$ facet reflections of an $n$-dimensional right-corner orthoscheme
$O(a_1,\dots,a_n)$ generate a group $W(a)\le\Isom(\R^n)$, and $W(a)$ is a quotient of the
right-angled Coxeter group $W_n$ on the orthogonality graph of the facets. This paper
determines what the quotient can do to the growth.

Three statements are proved unconditionally. First, a dimension-free dichotomy: if $Q=G/N$
is a quotient of a marked group and $K(N)$ is the length of a shortest nontrivial element of
$N$, then the sphere sizes of $Q$ and of $G$ agree in every degree below $\lceil K(N)/2\rceil$
and $Q$ is strictly smaller in that degree (Theorem~\ref{thm:dichotomy}). Second, the growth
series of $W_n$ is rational, equal to $(1+t)^{\lfloor n/2\rfloor+1}/K_{n+1}(t)$ for an explicit
integer polynomial $K_{n+1}$, with growth rate $r_n=1+2\cos\bigl(2\pi/(n+3)\bigr)$
(Theorem~\ref{thm:envelope}). Third, a complete classification of the rank-two relations:
$\ker(W_n\to W(a))$ meets a rank-two standard parabolic nontrivially if and only if the leg
sequence has two equal legs at an endpoint of the staircase or three consecutive equal legs
(Theorem~\ref{thm:rank2}). The Diophantine input is that three quartics in the flanking legs
have only trivial rational points, their Jacobians being the rank-zero curves \texttt{24a1},
\texttt{72a2} and $y^2=x^3-x$.

Combining the first two, the collision depth
$\cd_n(a)=\inf\{d:u_d(a)<[t^d]\,W_n(t)\}$ is exactly $K(\ker\rho_a)/2$. We show
$K(\ker\rho_a)\ge6$ for every positive leg tuple, so $\cd_n\ge3$ always, and we determine
$\cd_n$ from the two positional statistics $(e,t)$ of the leg sequence: unconditionally when
three consecutive legs are equal, and otherwise up to one explicitly stated hypothesis
(Theorem~\ref{thm:cd-general}).

What remains open is whether $\rho_a$ is injective at all for $n\ge3$; it is open in both
directions, and we state it as one problem in two forms (Problem~\ref{prob:alternative}). In
the plane $\rho_a$ is not injective, by an amenability argument, and that argument does not
extend: $O(n)$ contains a free subgroup of rank two for every $n\ge3$, so $\Isom(\R^n)$ is not
amenable as an abstract group (Remark~\ref{rem:amenability-fails}). We also record that
nonvanishing of the Schur-complement determinant $\det Q_n=-\prod_i a_i^2$
(Lemma~\ref{lem:uniform-detQ}) does not imply faithfulness, contrary to an argument used in an
earlier version of this material (Remark~\ref{rem:detQ-nonsequitur}).
\end{abstract}

\section{Introduction}\label{sec:intro}

An \emph{$n$-dimensional right-corner orthoscheme}, or path simplex,
$O=O(a_1,\dots,a_n)\subset\R^n$ is the simplex with vertices
\[
V_0=0,\qquad V_k=V_{k-1}+a_k e_k\quad(1\le k\le n),
\]
where $e_1,\dots,e_n$ is the standard basis and $a=(a_1,\dots,a_n)$ is a tuple of positive
reals, the \emph{legs}. Its $n+1$ facets carry reflections
$\widehat R_0,\dots,\widehat R_n\in\Isom(\R^n)$, and we write
$W(a)=\langle\widehat R_0,\dots,\widehat R_n\rangle$ for the group they generate.

Two facets $\widehat R_i,\widehat R_j$ with $|i-j|\ge2$ are perpendicular for every leg tuple
(Lemma~\ref{lem:normals}), while consecutive facets meet at an angle that depends on the legs
and is generically not a rational multiple of $\pi$. The largest group compatible with the
leg-independent relations is therefore the right-angled Coxeter group
\begin{equation}\label{eq:Wn-presentation}
W_n=\bigl\langle R_0,\dots,R_n\ \bigm|\ R_i^2=1,\ (R_iR_j)^2=1\ (|i-j|\ge2)\bigr\rangle,
\end{equation}
on the \emph{orthogonality graph} $\Gamma_n$ (vertices $0,\dots,n$, an edge whenever
$|i-j|\ge2$), and there is a surjection $\rho_a\colon W_n\twoheadrightarrow W(a)$ for every
$a$. Distinguishing $W_n$ from $W(a)$ is essential: writing $W_n$ for the geometric group would
make ``$\rho_a$ is faithful'' vacuous, and it is exactly the possible failure of injectivity
that this paper is about.

\subsection{Results}

The sphere sizes $u_d(a)=\#\{w\in W(a):\ell(w)=d\}$ are bounded above by the coefficients of
the growth series $W_n(t)$ of the envelope~\eqref{eq:Wn-presentation}, and the whole question is
where and by how much the bound fails to be sharp. Three unconditional theorems answer the
``where''.

\begin{itemize}[leftmargin=1.4em,itemsep=0.35em]
\item \textbf{The dichotomy} (Theorem~\ref{thm:dichotomy}). For any group $G$ with a finite
symmetric generating set and any $N\trianglelefteq G$ with $N\ne\{1\}$, put
$K(N)=\min\{\ell_G(g):g\in N\setminus\{1\}\}$. Then the sphere sizes of $G/N$ agree with those
of $G$ in every degree $d<\lceil K(N)/2\rceil$, and are strictly smaller in degree
$\lceil K(N)/2\rceil$. So a family of quotients of a fixed $G$ has a universal initial
segment, and the horizon of universality is half the shortest added relator.

\item \textbf{The envelope} (Theorem~\ref{thm:envelope}). $W_n(t)$ is rational, and is
$(1+t)^{\lfloor n/2\rfloor+1}/K_{n+1}(t)$ for the explicit integer polynomials $K_m$ of
\eqref{eq:Kdef}. Its growth rate is $r_n=1+2\cos\bigl(2\pi/(n+3)\bigr)$.

\item \textbf{The rank-two classification} (Theorem~\ref{thm:rank2}). For $a\in\Q_{>0}^n$, the
kernel of $\rho_a$ meets a rank-two standard parabolic $\langle R_i,R_{i+1}\rangle$
nontrivially if and only if $e(a)+t(a)\ge1$, where $e(a)$ counts equal legs at the two
endpoints of the staircase and $t(a)$ counts three consecutive equal legs. The shortest such
element has length $6$ when $t(a)\ge1$ and length $8$ when $t(a)=0<e(a)$.
\end{itemize}

The first two combine into the statement that gives the paper its shape. Define the
\emph{collision depth}
\[
\cd_n(a)=\inf\bigl\{d\ge0:\ u_d(a)<[t^d]\,W_n(t)\bigr\}\in\{1,2,\dots,\infty\}.
\]
By the dichotomy, $\cd_n(a)=\lceil K(\ker\rho_a)/2\rceil$, and since a relator of a reflection
group has even length this is $K(\ker\rho_a)/2$ (Corollary~\ref{cor:cd-is-K}). A short geometric
argument then shows $K(\ker\rho_a)\ge6$ for every positive leg tuple
(Lemma~\ref{lem:no-length-four}), so $\cd_n\ge3$ unconditionally, and the rank-two
classification pins the value: $\cd_n=3$ exactly when three consecutive legs are equal, and
$\cd_n\in\{3,4\}$ with the value $4$ in every computed case when only an endpoint pair is equal
(Theorem~\ref{thm:cd-general}). The remaining case $e=t=0$ is $\cd_n=\infty$ precisely when
$\rho_a$ is injective.

Whether $\rho_a$ is ever injective for $n\ge3$ is open in both directions. Section~\ref{sec:open}
states it as a single problem in two forms, generic and arithmetic, together with the two
partial results that bound it: the rank-two half is proved
(Theorem~\ref{thm:rank2}), and the one soft argument that settles the plane provably does not
extend (Remark~\ref{rem:amenability-fails}).

\subsection{What is and is not new}

The dichotomy is elementary. It is recorded here with a proof because it is what turns the two
halves of this paper into statements about each other, and because its deviation half needs a
case distinction on the parity of $K(N)$ that is easy to omit; but it is not a difficult
argument, and the phenomenon it describes, that a family of quotients of a common group shares
an initial segment of its growth, is a general one of which the orthoscheme family is an
instance. The present paper claims the instance, not the principle.

Both ingredients of $r_n$ are classical and are used as such. Steinberg's formula for the
Poincar\'e series of a Coxeter group is standard \cite{Steinberg,humphreys90,Bourbaki}, growth
series of right-angled Coxeter groups are treated systematically by Kellerhals and
Perren~\cite{kellerhals-perren}, and the independence polynomials of paths, their Binet formula,
and the location of their roots are classical facts about the Fibonacci--Pell family. What is
asserted is the identification: that the growth rate of the $n$-dimensional right-corner
orthoscheme envelope is the largest of those roots at $m=n+3$. We have not located that
identification in the literature and make no claim of novelty for the ingredients. A reader who
knows the Coxeter growth-series literature should read Theorem~\ref{thm:envelope} as an assembly
of known pieces.

The Diophantine part of Theorem~\ref{thm:rank2} rests on three rank-zero Mordell--Weil
computations that are in the standard tables \cite{cremona-1997,lmfdb-24a1,lmfdb-72a2}; what is
done here is the reduction of the dihedral condition to those three curves.

Two points of contact should be recorded. The all-equal tuple $a_1=\cdots=a_n$ is not a generic
point of this family: there the orthoscheme group is the affine Weyl group $\widetilde C_n$,
whose growth is classical and whose Poincar\'e series is given by Bott's formula, and which
reproduces the $\cd_n=3$ rows of Table~\ref{tab:atlas}. In the planar case $n=2$, Isola and
Piergallini~\cite{isola-piergallini} prove that the generic triangle reflection group is not
finitely presented, which already implies that $\rho_a$ is not injective there; that is a second
reason, besides the amenability argument of Proposition~\ref{prop:plane}, why the plane is not a
guide to $n\ge3$.

\subsection{Companion papers}

The planar case $n=2$ is treated separately \cite{Bonfioli1}: there the metric has a
lamplighter-type closed form, the deviation from the envelope has been computed past its onset,
and the resulting growth series (\textup{A396406}, growth rate $\beta_2=1.4916\ldots$) is
transcendental over $\Q(x)$. Nothing below is quoted from that work except the two numerical
facts recorded in Corollary~\ref{cor:plane}, and the transcendence questions, which are special
to $n=2$, are not attempted here. Non-right-corner orthoschemes are outside the scope of this
paper.

\section{Orthoschemes and their Coxeter envelope}\label{sec:setting}

Throughout, $n\ge2$, and
\[
U=\{a=(a_1,\dots,a_n)\in\R^n:\ a_i>0\ \text{for all }i\}
\]
is the leg space. Everything below is invariant under simultaneous scaling of the legs, so one
may normalise if convenient; we do not.

\subsection{Facet normals}

\begin{lemma}[Facet normals]\label{lem:normals}
With $V_0=0$ and $V_k=\sum_{i\le k}a_ie_i$, the facet of $O(a)$ opposite $V_j$ has normal
\[
m_0=e_1,\qquad m_j=a_{j+1}e_j-a_je_{j+1}\ \ (1\le j\le n-1),\qquad m_n=e_n .
\]
Consequently $\operatorname{supp}(m_0)=\{1\}$, $\operatorname{supp}(m_j)=\{j,j+1\}$ and
$\operatorname{supp}(m_n)=\{n\}$, so $m_i\cdot m_j=0$ whenever $|i-j|\ge2$, while
$m_i\cdot m_{i+1}\ne0$ for every $i$ and every $a\in U$.
\end{lemma}

\begin{proof}
The facet opposite $V_j$ is the affine hull of $\{V_i:i\ne j\}$. For $j=0$ the difference
vectors $V_{k+1}-V_k=a_{k+1}e_{k+1}$, $1\le k\le n-1$, span $\langle e_2,\dots,e_n\rangle$,
whence $m_0=e_1$. For $j=n$ the vertices $V_0,\dots,V_{n-1}$ span
$\langle e_1,\dots,e_{n-1}\rangle$, whence $m_n=e_n$. For $1\le j\le n-1$ the difference
vectors are $a_ke_k$ for $k\le j-1$, then $V_{j+1}-V_{j-1}=a_je_j+a_{j+1}e_{j+1}$, then
$a_ke_k$ for $k\ge j+2$. Orthogonality to the first and third families confines $m_j$ to
coordinates $\{j,j+1\}$, and orthogonality to the middle vector forces
$m_j\propto a_{j+1}e_j-a_je_{j+1}$. Two supports of the listed form are disjoint exactly when
$|i-j|\ge2$. For adjacent indices, $m_0\cdot m_1=a_2$, $m_j\cdot m_{j+1}=-a_ja_{j+2}$ for
$1\le j\le n-2$, and $m_{n-1}\cdot m_n=-a_{n-1}$, all nonzero because every $a_i>0$.
\end{proof}

Lemma~\ref{lem:normals} holds for all legs, with no genericity hypothesis and no case check.
Hence $\widehat R_i^2=1$ always and $\widehat R_i\widehat R_j=\widehat R_j\widehat R_i$ whenever
$|i-j|\ge2$, so the assignment $R_i\mapsto\widehat R_i$ defines a surjection
\[
\rho_a\colon W_n\twoheadrightarrow W(a)\le\Isom(\R^n)
\]
from the right-angled Coxeter group~\eqref{eq:Wn-presentation} on the orthogonality graph
$\Gamma_n$ \cite{Davis}. The abstract group $W_n$ does not depend on $a$.

\subsection{Counting conventions}

For a group $H$ with a marked finite generating set we write $\ell_H$ for the word length and
$s_d(H)=\#\{h\in H:\ell_H(h)=d\}$ for the sphere sizes. We abbreviate
\[
u_d(a)=s_d\bigl(W(a)\bigr),\qquad W_n(t)=\sum_{d\ge0}s_d(W_n)\,t^d ,
\]
the generators being $\widehat R_0,\dots,\widehat R_n$ and $R_0,\dots,R_n$ respectively. Since
$O(a)$ is a fundamental domain for the action of $W(a)$, $u_d(a)$ is also the number of distinct
chambers $g\cdot O(a)$ at word distance $d$, which is what the breadth-first search of
\S\ref{sec:cert} enumerates. Being the sphere sizes of a quotient of $W_n$, they satisfy
$u_d(a)\le[t^d]\,W_n(t)$ for every $d$.

Finally, set
\[
K(a)=\min\bigl\{\ell_{W_n}(w):w\in\ker\rho_a\setminus\{1\}\bigr\}\in\{1,2,\dots\}\cup\{\infty\},
\]
with $K(a)=\infty$ when $\rho_a$ is injective.

\subsection{Generic rigidity}

Before computing anything we record that ``the generic orthoscheme group'' is a well-defined
object: it does not depend on which generic shape one picks. Write
$N_{\mathrm{gen}}=\bigcap_{a\in U}\ker\rho_a\trianglelefteq W_n$.

\begin{theorem}[Generic rigidity]\label{thm:rigidity}
For every $n\ge2$ there is a co-null subset $U_{\mathrm{gen}}\subseteq U$, the complement of a
countable union of proper real-algebraic subvarieties, such that
$\ker\rho_a=N_{\mathrm{gen}}$, hence $W(a)\cong W_n/N_{\mathrm{gen}}$, for every
$a\in U_{\mathrm{gen}}$. Consequently $u_d(a)$ is constant on $U_{\mathrm{gen}}$, equal to
$s_d(W_n/N_{\mathrm{gen}})$, and this generic value is the maximum of $u_d$ over all of $U$.
Moreover $U_{\mathrm{gen}}$ contains every $a$ whose coordinates are algebraically independent
over $\Q$; in particular it is dense.
\end{theorem}

\begin{proof}
Each facet normal $m_j(a)$ is a fixed polynomial vector in $a$ by Lemma~\ref{lem:normals}, so
the reflection $\widehat R_j=I-2\,m_jm_j^{\!\top}/(m_j\!\cdot\!m_j)$, together with its
translation part, has entries in $\Q(a)$ with denominator the nowhere-vanishing polynomial
$m_j\!\cdot\!m_j$. Hence for each $w\in W_n$, choosing any word representative, every entry of
$\rho_a(w)$ lies in $\Q(a)$, independently of the representative because $\rho_a$ is well
defined. The locus $V_w=\{a\in U:\rho_a(w)=I\}$ is the simultaneous vanishing of finitely many
rational functions cleared of their nonvanishing denominators, hence Zariski-closed in the
irreducible variety $U$. If $w\notin N_{\mathrm{gen}}$ then $\rho_a(w)\ne I$ for some $a$, so
$V_w$ is a proper closed subvariety, of measure zero. As $W_n$ is countable,
$B:=\bigcup_{w\in W_n\setminus N_{\mathrm{gen}}}V_w$ is a countable union of proper
subvarieties, and $U_{\mathrm{gen}}:=U\setminus B$ is co-null and dense.

For $a\in U_{\mathrm{gen}}$ one has $N_{\mathrm{gen}}\subseteq\ker\rho_a$ by definition, and
$\ker\rho_a\subseteq N_{\mathrm{gen}}$ because any
$w\in\ker\rho_a\setminus N_{\mathrm{gen}}$ would place $a$ in $V_w\subseteq B$. Isomorphic
groups with marked generators have identical sphere sizes, giving the constant value; for
arbitrary $a$, $\ker\rho_a\supseteq N_{\mathrm{gen}}$ makes $W(a)$ a quotient of
$W_n/N_{\mathrm{gen}}$, so $u_d(a)\le s_d(W_n/N_{\mathrm{gen}})$. Finally each $V_w$ is defined
over $\Q$; if the coordinates of $a$ are algebraically independent over $\Q$ then $a$ lies on no
proper $\Q$-subvariety, so $a\notin B$. Such $a$ exist since $\dim U=n\ge2$.
\end{proof}

\begin{corollary}[One faithful shape suffices]\label{cor:one-shape}
If $\rho_a$ is injective for a single $a\in U$, then $N_{\mathrm{gen}}=\{1\}$ and $\rho_b$ is
injective for every $b\in U_{\mathrm{gen}}$. Conversely, if $\rho_b$ fails to be injective for
some $b\in U_{\mathrm{gen}}$, then $\rho_a$ fails to be injective for every $a\in U$.
\end{corollary}

\begin{proof}
$N_{\mathrm{gen}}\subseteq\ker\rho_a=\{1\}$, and $\ker\rho_b=N_{\mathrm{gen}}$ for
$b\in U_{\mathrm{gen}}$ by Theorem~\ref{thm:rigidity}. The converse is the contrapositive,
using $N_{\mathrm{gen}}\subseteq\ker\rho_a$ for every $a$.
\end{proof}

Corollary~\ref{cor:one-shape} is what allows the generic question and the arithmetic question of
\S\ref{sec:open} to be treated as two forms of one problem in one direction. It gives nothing in
the other direction: a specific rational tuple need not lie in $U_{\mathrm{gen}}$, and
$U_{\mathrm{gen}}$ contains no tuple with algebraically dependent coordinates, hence no rational
tuple at all.

\section{The universality--deviation dichotomy}\label{sec:dichotomy}

\begin{theorem}[Universality--deviation dichotomy]\label{thm:dichotomy}
Let $G=\langle S\rangle$ be a group with finite symmetric generating set $S$, let
$N\trianglelefteq G$ with $N\ne\{1\}$, write $Q=G/N$ and
$K=K(N)=\min\{\ell_G(g):g\in N\setminus\{1\}\}$, and put $h=\lceil K/2\rceil$. Then
$s_d(Q)=s_d(G)$ for every $d<h$, and $s_h(Q)<s_h(G)$.
\end{theorem}

\begin{proof}
Write $\pi:G\to Q$.

\emph{Agreement below $h$.} Let $d<h$. If $d\ge1$ then $2d\le2h-2\le K-1<K$, since $2h\le K+1$.
Suppose $g,g'\in B_d(G)$ satisfy $\pi g=\pi g'$. Then $g(g')^{-1}\in N$ and
$\ell_G(g(g')^{-1})\le2d<K$, so $g=g'$; thus $\pi$ is injective on $B_d(G)$. It is also
length-preserving there: if $\ell_G(g)=d$ and $\ell_Q(\pi g)=e<d$, pick $g'$ with
$\ell_G(g')=e$ and $\pi g'=\pi g$; then $\ell_G(g(g')^{-1})\le d+e<2d<K$ forces $g=g'$,
contradicting $e<d$. Finally $\pi$ maps the sphere of radius $d$ in $G$ onto the sphere of
radius $d$ in $Q$, because $\ell_Q(q)=\min\{\ell_G(g):\pi g=q\}$ realises every $q$ with
$\ell_Q(q)=d$ by some $g$ of length exactly $d$. Hence $s_d(Q)=s_d(G)$.

\emph{Strict drop at $h$.} Choose $g_0\in N\setminus\{1\}$ with $\ell_G(g_0)=K$ and split a
geodesic word for it as $g_0=w_1w_2$ with $\ell_G(w_1)=h$ and
$\ell_G(w_2)=K-h=\lfloor K/2\rfloor$; both factors are geodesic as written. From $w_1w_2\in N$
we get $\pi(w_1)=\pi(w_2^{-1})$, and $w_1\ne w_2^{-1}$ since otherwise $g_0=1$. Surjectivity of
$\pi$ from the radius-$h$ sphere of $G$ onto that of $Q$, established in the previous paragraph
for $d<h$, holds verbatim for $d=h$.

If $K$ is even then $h=K/2=\lfloor K/2\rfloor$, so $w_1$ and $w_2^{-1}$ are two distinct
elements of the sphere of radius $h$ in $G$ with the same image, and $s_h(Q)\le s_h(G)-1$.

If $K$ is odd then $h=(K+1)/2$ and $\ell_G(w_2^{-1})=h-1$, so $\ell_Q(\pi w_1)\le h-1<h$ and the
image of $w_1$ does not lie on the sphere of radius $h$ in $Q$. That sphere is therefore covered
by $\pi$ applied to the radius-$h$ sphere of $G$ with $w_1$ removed, and again
$s_h(Q)\le s_h(G)-1$.
\end{proof}

The parity split at the end is not cosmetic. When $K$ is odd the two halves of a shortest
relator lie on spheres of different radii, so no two points of the upstairs sphere need collide,
and the count drops for the other reason. Both cases are machine-checked
(\S\ref{sec:cert}).

Applied to $\rho_a$, Theorem~\ref{thm:dichotomy} gives an exact alternative with no genericity
and no dimension hypothesis.

\begin{corollary}[The envelope alternative]\label{cor:alternative}
Exactly one of the following holds.
\begin{enumerate}[label=(\roman*),leftmargin=2.2em]
\item $\rho_a$ is injective. Then $W(a)\cong W_n$ and $u_d(a)=[t^d]\,W_n(t)$ for every $d$; in
particular the growth series of $W(a)$ is the rational function of
Theorem~\ref{thm:envelope}, with growth rate $r_n$.
\item $\rho_a$ is not injective. Then $K(a)<\infty$, $u_d(a)=[t^d]\,W_n(t)$ for every
$d<\lceil K(a)/2\rceil$, and $u_d(a)<[t^d]\,W_n(t)$ at $d=\lceil K(a)/2\rceil$.
\end{enumerate}
\end{corollary}

\begin{proof}
If $\ker\rho_a=\{1\}$ then $\rho_a$ is an isomorphism of marked groups and (i) is immediate.
Otherwise $K(a)$ is finite and Theorem~\ref{thm:dichotomy} with $G=W_n$, $N=\ker\rho_a$ gives
(ii).
\end{proof}

\begin{corollary}[The collision depth is half the shortest relator]\label{cor:cd-is-K}
For every $a\in U$, $\ \cd_n(a)=\lceil K(a)/2\rceil$, and $K(a)$ is even, so
$\cd_n(a)=K(a)/2$.
\end{corollary}

\begin{proof}
The first equality is Corollary~\ref{cor:alternative}: below $\lceil K(a)/2\rceil$ the two
sequences agree and at $\lceil K(a)/2\rceil$ they differ. For parity, let
$\varepsilon:W_n\to\{\pm1\}$ be the sign character, $\varepsilon(R_i)=-1$, which is well defined
because all defining relators of \eqref{eq:Wn-presentation} have even length; then
$\varepsilon(w)=(-1)^{\ell(w)}$. On the other side, $\det\circ\rho_a$ is a homomorphism
$W_n\to\{\pm1\}$ sending each $R_i$ to $-1$, so $\det\rho_a(w)=\varepsilon(w)$. If
$w\in\ker\rho_a$ then $\det\rho_a(w)=1$, hence $\ell(w)$ is even.
\end{proof}

\begin{lemma}[No relator of length four]\label{lem:no-length-four}
For every $n\ge2$ and every $a\in U$, $K(a)\ge6$. Consequently $\cd_n(a)\ge3$.
\end{lemma}

\begin{proof}
By Corollary~\ref{cor:cd-is-K}, $K(a)$ is even, so it suffices to exclude $K(a)\in\{2,4\}$.

Write $H_i$ for the facet hyperplane carrying $\widehat R_i$. The $H_i$ are the $n+1$ facet
hyperplanes of a nondegenerate simplex, so they are pairwise distinct and pairwise
non-parallel, and $H_i\cap H_j$ is the affine hull of $\{V_m:m\ne i,j\}$, of dimension $n-2$.
Distinct index pairs give distinct intersections, because an $(n-2)$-flat contains exactly
$n-1$ of the $n+1$ affinely independent vertices.

$K(a)=2$ would give a geodesic word $R_iR_j$ with $i\ne j$ and
$\widehat R_i=\widehat R_j$, impossible since $H_i\ne H_j$.

Suppose $K(a)=4$, realised by a geodesic word $R_iR_jR_kR_l$ with $i\ne j$, $j\ne k$, $k\ne l$
and $\widehat R_i\widehat R_j\widehat R_k\widehat R_l=1$, that is
$\widehat R_i\widehat R_j=\widehat R_l\widehat R_k$. For $H_i\ne H_j$ non-parallel, the isometry
$\widehat R_i\widehat R_j$ is a rotation about $H_i\cap H_j$ through twice the angle between the
hyperplanes, an angle in $(0,2\pi)$; its fixed-point set is therefore exactly $H_i\cap H_j$. The
same applies to $\widehat R_l\widehat R_k$, so $H_i\cap H_j=H_k\cap H_l$ and hence
$\{i,j\}=\{k,l\}$ by the previous paragraph. Since $j\ne k$ this forces $i=k$ and $j=l$, so the
relation reads $\widehat R_i\widehat R_j=\widehat R_j\widehat R_i$, that is $H_i\perp H_j$. For
$|i-j|\ge2$ that relation already holds in $W_n$ and the word $R_iR_jR_iR_j$ is not geodesic;
for $|i-j|=1$ perpendicularity contradicts $m_i\cdot m_{i+1}\ne0$ in
Lemma~\ref{lem:normals}. Both branches are impossible, so $K(a)\ne4$.
\end{proof}

\section{The envelope in closed form}\label{sec:envelope}

\begin{theorem}[The envelope is rational]\label{thm:envelope}
Let $c_{\Gamma_n}(x)=\sum_{C\ \mathrm{clique}}x^{|C|}$ be the clique polynomial of $\Gamma_n$.
The growth series of $W_n$ is rational and satisfies
\[
\frac{1}{W_n(t)}=c_{\Gamma_n}\!\Bigl(\frac{-t}{1+t}\Bigr).
\]
Equivalently $W_n(t)=(1+t)^{n+1}/J_{n+1}(t)$, where
\[
J_m=(1+t)\bigl(J_{m-1}-t\,J_{m-2}\bigr),\qquad J_0=J_1=1 .
\]
One has $J_m=(1+t)^{\lfloor m/2\rfloor}K_m$ with
\begin{equation}\label{eq:Kdef}
K_0=K_1=1,\qquad K_{2j}=K_{2j-1}-tK_{2j-2},\qquad K_{2j+1}=(1+t)K_{2j}-tK_{2j-1},
\end{equation}
so that
\[
W_n(t)=\frac{(1+t)^{\lfloor n/2\rfloor+1}}{K_{n+1}(t)} .
\]
The reciprocals of the roots of $K_{n+1}$ in $(1/3,1)$ are the numbers
$1+2\cos\frac{2k\pi}{n+3}$, and the growth rate is the largest of them,
\[
r_n=1+2\cos\frac{2\pi}{n+3}.
\]
\end{theorem}

\begin{proof}
$W_n$ is a right-angled Coxeter group, so its finite standard parabolics are exactly the cliques
of $\Gamma_n$, each a group $(\Z/2)^{|C|}$ with growth polynomial $(1+t)^{|C|}$ \cite{Davis}.
Steinberg's formula \cite{Steinberg,Bourbaki,humphreys90} gives
\[
\frac{1}{W_n(t^{-1})}=\sum_{C\ \mathrm{clique}}\frac{(-1)^{|C|}}{(1+t)^{|C|}}
=c_{\Gamma_n}\!\Bigl(\frac{-1}{1+t}\Bigr).
\]
Substituting $t\mapsto t^{-1}$ and simplifying $-1/(1+t^{-1})=-t/(1+t)$ yields the displayed
identity for $1/W_n(t)$. The substitution must be carried inside the clique polynomial;
retaining the argument $-1/(1+t)$ gives a different rational function, already at $n=2$.

The cliques of $\Gamma_n$ are the independent sets of its complement, the path $P_{n+1}$, so
$c_{\Gamma_n}=I_{n+1}$, the independence polynomial of the path on $n+1$ vertices, which obeys
$I_m=I_{m-1}+xI_{m-2}$, $I_0=1$, $I_1=1+x$, and has the coefficient form
$I_m(x)=\sum_k\binom{m+1-k}{k}x^k$. Put $J_m(t)=(1+t)^mI_m(-t/(1+t))$. Multiplying the
recurrence by $(1+t)^m$ gives $J_m=(1+t)J_{m-1}-t(1+t)J_{m-2}$, which is the stated recurrence,
with $J_0=1$ and $J_1=(1+t)(1-t/(1+t))=1$. By construction
$1/W_n(t)=I_{n+1}(-t/(1+t))=J_{n+1}(t)/(1+t)^{n+1}$. The factorisation
$J_m=(1+t)^{\lfloor m/2\rfloor}K_m$ and the two-step recurrences \eqref{eq:Kdef} follow by
induction on $m$, splitting on the parity of $m$; cancelling $(1+t)^{\lfloor(n+1)/2\rfloor}$
against $(1+t)^{n+1}$ leaves the exponent $\lceil(n+1)/2\rceil=\lfloor n/2\rfloor+1$.

For the roots, fix $t>0$ and solve the linear recurrence. Its characteristic equation is
$\lambda^2-(1+t)\lambda+t(1+t)=0$, of discriminant $(1+t)(1-3t)$.

If $0<t<1/3$ both roots are real with $0<\lambda_-<\lambda_+<1$, their sum being $1+t\le4/3$ and
their product $t(1+t)<4/9$. Writing $J_m=\alpha\lambda_+^m+\beta\lambda_-^m$, the initial
conditions give $\alpha=(1-\lambda_-)/(\lambda_+-\lambda_-)>0$ and
$\beta=(\lambda_+-1)/(\lambda_+-\lambda_-)<0$ with $\alpha+\beta=1$. Then
$J_m=\lambda_-^m(\alpha(\lambda_+/\lambda_-)^m+\beta)$, and the bracket is increasing in $m$
with value $1$ at $m=0$, so $J_m>0$. At $t=1/3$ the root is double and $J_m=(1+m/2)(2/3)^m>0$.
Hence $J_m$ has no root in $(0,1/3]$.

If $t>1/3$ the roots are $\lambda_\pm=\varrho e^{\pm i\theta}$ with $\varrho=\sqrt{t(1+t)}$ and
$\cos\theta=\sqrt{(1+t)/(4t)}$, so $\theta$ ranges over $(0,\pi/3)$ and is strictly increasing
in $t$, because $(1+t)/(4t)=\tfrac14+\tfrac1{4t}$ is strictly decreasing. Then
$J_m=\varrho^m(\cos m\theta-c\sin m\theta)$, where $J_0=1$ is automatic and $J_1=1$ forces
$c\varrho\sin\theta=(t-1)/2$, that is $c=(t-1)/\sqrt{(1+t)(3t-1)}$. Since
$\sin^2\theta=(3t-1)/(4t)$, a direct computation gives
$\cot2\theta=(1-t)/\sqrt{(1+t)(3t-1)}=-c$. Therefore
\[
J_m=0\iff\cot m\theta=c=-\cot2\theta
\iff\frac{\sin\bigl((m+2)\theta\bigr)}{\sin m\theta\,\sin2\theta}=0
\iff(m+2)\theta\in\pi\Z .
\]
With $m=n+1$ the roots are exactly $\theta=k\pi/(n+3)$, $1\le k<(n+3)/3$. Inverting
$\cos^2\theta=(1+t)/(4t)$ gives $t=1/(4\cos^2\theta-1)=1/(1+2\cos2\theta)$, so the corresponding
$t$-values are $1/\bigl(1+2\cos\frac{2k\pi}{n+3}\bigr)$. As $\theta$ increases with $t$, the
smallest positive root is $k=1$; it exceeds $1/3$ because $1+2\cos\frac{2\pi}{n+3}<3$, and
$1+t>0$ there, so it is a root of $K_{n+1}$ and not of the stripped factor. Hence
$r_n=1+2\cos\frac{2\pi}{n+3}$.
\end{proof}

\begin{remark}[Not a path spectral radius]\label{rem:notpath}
The number $2\cos\frac{2\pi}{n+3}$ is not the spectral radius of the path $P_{n+1}$, which is
$2\cos\frac{\pi}{n+2}$; the two disagree for every $n\ge2$ (at $n=2$ they are $0.618\ldots$ and
$1.414\ldots$). The angles produced by the proof are $k\pi/(n+3)$, and the growth rate uses
$\cos2\theta$ at $\theta=\pi/(n+3)$, not $\cos\theta$.
\end{remark}

\begin{remark}[The clique polynomial is a Fibonacci--Pell polynomial]\label{rem:pell}
The coefficient form $I_{n+1}(x)=\sum_k\binom{n+2-k}{k}x^k$ used above is the statement that the
clique count of $\Gamma_n$ in size $k$ is $\binom{n+2-k}{k}$, and it identifies
$c_{\Gamma_n}$ with the Fibonacci--Pell polynomial $P_{n+3}$, where $P_0=0$, $P_1=1$ and
$P_m=P_{m-1}+yP_{m-2}$: the two families satisfy the same three-term recursion, by Pascal's
identity $\binom{n-k+2}{k}=\binom{n-k+1}{k}+\binom{n-k+1}{k-1}$, and agree at the bottom. This
gives a second reading of the roots. The characteristic equation $x^2=x+y$ has roots
$\alpha,\beta=(1\pm\sqrt{1+4y})/2$, the Binet formula yields
$P_m(y)=(\alpha^m-\beta^m)/(\alpha-\beta)$, so $P_m(y)=0$ iff $(\alpha/\beta)^m=1$; setting
$\alpha/\beta=e^{2\pi ij/m}$ gives $y_j=-1/\bigl(4\cos^2(\pi j/m)\bigr)$, and translating back
through $y=-1/(1+t)$ with $4\cos^2\vartheta-1=1+2\cos2\vartheta$ returns the same values
$1+2\cos\frac{2\pi j}{n+3}$. The Pascal step, the Binet formula and the translation identity are
machine-checked (\S\ref{sec:cert}); the induction that assembles the Pascal step into the
identification $c_{\Gamma_n}=P_{n+3}$ is a paper proof.
\end{remark}

\begin{corollary}[Asymptotic regime]\label{cor:asymptotic}
$r_n\to3$ as $n\to\infty$, with
\[
3-r_n=4\sin^2\!\Bigl(\frac{\pi}{n+3}\Bigr)\sim\frac{4\pi^2}{(n+3)^2}.
\]
\end{corollary}

Since $2\cos(2\pi/m)$ has degree $\phi(m)/2$ over $\Q$, where $\phi$ denotes Euler's totient,
the rate $r_n$ is rational or a quadratic surd exactly when $\phi(n+3)\le4$, that is
$n+3\in\{5,6,8,10,12\}$; see Table~\ref{tab:n-dim-rates}, where $\varphi$ denotes the golden
ratio.

\begin{table}[ht!]
\centering
\small
\renewcommand{\arraystretch}{1.2}
\begin{tabular}{c l r l}
\toprule
$n$ & $W_n(t)$ & $r_n$ & algebraic form \\
\midrule
2  & $(1+t)^2/(1-t-t^2)$                        & $1.6180\ldots$ & $\varphi=(1+\sqrt5)/2$ (golden) \\
3  & $(1+t)^2/(1-2t)$                           & $2.0000$       & exactly $2$ \\
4  & $(1+t)^3/(1-2t-t^2+t^3)$                   & $2.2470\ldots$ & $1+2\cos(2\pi/7)$, cubic \\
5  & $(1+t)^3/(1-3t+t^2+t^3)$                   & $2.4142\ldots$ & $1+\sqrt2$ (silver) \\
6  & $(1+t)^4/(1-3t+3t^3)$                      & $2.5321\ldots$ & $1+2\cos(2\pi/9)$, cubic \\
7  & $(1+t)^4/(1-4t+3t^2+2t^3-t^4)$             & $2.6180\ldots$ & $\varphi^2=(3+\sqrt5)/2$ \\
8  & $(1+t)^5/(1-4t+2t^2+5t^3-2t^4-t^5)$        & $2.6825\ldots$ & $1+2\cos(2\pi/11)$, degree $5$ \\
9  & $(1+t)^5/(1-5t+6t^2+2t^3-4t^4)$            & $2.7321\ldots$ & $1+\sqrt3$ \\
10 & $(1+t)^6/(1-5t+5t^2+6t^3-7t^4-2t^5+t^6)$   & $2.7709\ldots$ & $1+2\cos(2\pi/13)$, degree $6$ \\
\bottomrule
\end{tabular}
\caption{The envelope $W_n(t)=(1+t)^{\lfloor n/2\rfloor+1}/K_{n+1}(t)$ and its growth rate
$r_n=1+2\cos(2\pi/(n+3))$. The recurrence \eqref{eq:Kdef}, its agreement with the Steinberg
expansion, the exactness of the division by $(1+t)^{\lfloor(n+1)/2\rfloor}$, and the
identification of the smallest positive root of $K_{n+1}$ with $1/r_n$ are certified for
$2\le n\le30$ by the program of \S\ref{sec:cert}.}
\label{tab:n-dim-rates}
\end{table}

\begin{remark}[The orthoscheme sequences in the OEIS]\label{rem:oeis}
For pairwise distinct legs the layer counts $u_d(a)$ agree with the envelope as far as they have
been computed, and the resulting sequences are recorded in the On-Line Encyclopedia of Integer
Sequences: A397439 ($n=3$), A397437 ($n=4$), A396927 ($n=5$) and A397438 ($n=6$), with the $n=7$
entry in preparation. Each of these is the coefficient sequence of the rational function
$W_n(t)$ of Table~\ref{tab:n-dim-rates}, with dominant pole $1/r_n$. The planar entry A396406 is
of a different nature: it records the deviated planar sequence, which agrees with $W_2(t)$ only
up to degree $9$ and whose generating function is transcendental over $\Q(x)$
\cite{Bonfioli1}; see Corollary~\ref{cor:plane}.
\end{remark}

\section{Rank-two relations}\label{sec:rank2}

This section determines exactly when $\ker\rho_a$ meets a rank-two standard parabolic
$\langle R_i,R_{i+1}\rangle$. Throughout it, the legs are positive rationals. By
Lemma~\ref{lem:normals} the only pairs that can carry a nontrivial dihedral relation are the
consecutive ones, and their cosines are
\begin{equation}\label{eq:cosines}
\cos\theta_{0,1}=\frac{a_2}{\sqrt{a_1^2+a_2^2}},\qquad
\cos\theta_{n-1,n}=\frac{a_{n-1}}{\sqrt{a_{n-1}^2+a_n^2}},\qquad
\cos\theta_{i,i+1}=\frac{a_ia_{i+2}}{\sqrt{(a_i^2+a_{i+1}^2)(a_{i+1}^2+a_{i+2}^2)}}
\end{equation}
for $1\le i\le n-2$. In each case $\cos^2\theta\in\Q$. By the extended Niven theorem, a
consequence of Conway--Jones \cite{conway-jones} (if $\theta\in\Q\pi$ and $\cos^2\theta\in\Q$
then $\cos\theta\in\{0,\pm\tfrac12,\pm\tfrac{\sqrt2}2,\pm\tfrac{\sqrt3}2,\pm1\}$), a relation
$(R_iR_{i+1})^m=1$ in $W(a)$ forces
\[
\cos^2\theta_{i,i+1}\in\{0,\tfrac14,\tfrac12,\tfrac34,1\}.
\]

\subsection{The two positional statistics}

For $a\in\Q_{>0}^n$ set
\begin{align*}
e(a)&=\#\bigl\{i\in\{1,n-1\}:a_i=a_{i+1}\bigr\}\in\{0,1,2\},\\
t(a)&=\#\bigl\{i:1\le i\le n-2,\ a_i=a_{i+1}=a_{i+2}\bigr\}.
\end{align*}
Thus $e$ counts equal pairs at the two endpoint positions of the staircase and $t$ counts
interior triples of three consecutive equal legs. An equal pair sitting strictly inside the
staircase, such as $a_2=a_3$ in $(1,2,2,3)$ at $n=4$, contributes to neither.

\subsection{Boundary pairs}

\begin{lemma}[Boundary pairs]\label{lem:boundary}
A boundary pair carries an admissible value of $\cos^2\theta$ if and only if the two legs
involved are equal, in which case $\cos^2\theta=\tfrac12$ and $\theta=\pi/4$.
\end{lemma}

\begin{proof}
Take the pair $(0,1)$, so $\cos^2\theta=a_2^2/(a_1^2+a_2^2)$ by \eqref{eq:cosines}; the pair
$(n-1,n)$ is symmetric. The value $0$ forces $a_2=0$ and the value $1$ forces $a_1=0$, both
excluded by positivity. The value $\tfrac14$ gives $4a_2^2=a_1^2+a_2^2$, that is
$(a_1/a_2)^2=3$, and the value $\tfrac34$ gives $(a_2/a_1)^2=3$; neither is possible over $\Q$.
The value $\tfrac12$ gives $a_1^2=a_2^2$, that is $a_1=a_2$, and conversely $a_1=a_2$ gives
$\cos\theta=1/\sqrt2$.
\end{proof}

All five cases of Lemma~\ref{lem:boundary} are formalised in Lean~4 (\S\ref{sec:cert}).

\subsection{Interior pairs: three Diophantine quartics}

At an interior pair write $x=a_i$, $y=a_{i+1}$, $z=a_{i+2}$, so that by \eqref{eq:cosines}
\[
\cos^2\theta_{i,i+1}=\frac{x^2z^2}{(x^2+y^2)(y^2+z^2)} .
\]
Requiring this to equal a constant $c$ and setting $u=y^2$ gives the quadratic
\begin{equation}\label{eq:uquad}
c\,u^2+c\,(x^2+z^2)\,u+(c-1)\,x^2z^2=0 ,
\end{equation}
which has a rational root only if its discriminant
$\Delta=c^2(x^2+z^2)^2-4c(c-1)x^2z^2$ is the square of a rational. Substituting the three
interior values of $c$ and clearing the constant factors gives, in the two flanking legs
$x=a_i$ and $z=a_{i+2}$, the Diophantine \emph{quartics}
\[
\mathcal{V}_{1/4}:\ x^4+14x^2z^2+z^4=k^2,\qquad
\mathcal{V}_{1/2}:\ x^4+6x^2z^2+z^4=k^2,\qquad
\mathcal{V}_{3/4}:\ 9x^4+30x^2z^2+9z^4=k^2 .
\]
It matters that these are quartics in the legs and not conics in $X=x^2$, $Z=z^2$. Written in
$X$ and $Z$ the three equations define conics with infinitely many rational points, among them
$(X,Z)=(3,2)$ on $X^2+6XZ+Z^2=k^2$ with $k=7$, and $(X,Z)=(2,5)$ and $(5,9)$ on the other two.
Those points are irrelevant here precisely because $X$ and $Z$ must themselves be squares of
legs.

\begin{lemma}\label{lem:cremona-empty}
Each of $\mathcal{V}_{1/4}$, $\mathcal{V}_{1/2}$, $\mathcal{V}_{3/4}$ has only trivial solutions
over $\Q_{>0}^2$ in the legs $(x,z)$: every rational point has $xz=0$ or $x=z$.
\end{lemma}

\begin{proof}
\emph{$\mathcal{V}_{1/2}$.} The quartic $x^4+6x^2z^2+z^4=k^2$ is birationally equivalent to the
elliptic curve $E_{1/2}:y^2=x^3-x$ (Mordell~\cite{mordell-1969}, \S17). $E_{1/2}$ has
Mordell--Weil rank $0$ over $\Q$ by Fermat's classical descent (Cassels~\cite{cassels-1991},
Ch.~II), and its full rational-point set is the $2$-torsion $\{\infty,(0,0),(\pm1,0)\}$, pulling
back to $xz=0$ or $x=z$.

\emph{$\mathcal{V}_{1/4}$.} The Jacobian elliptic curve is $E_{1/4}:y^2=x(x-12)(x-16)$, which is
Cremona \texttt{24a1} (conductor $24$, $j=35\,152/9$, Mordell--Weil rank $0$ over $\Q$, torsion
$\Z/4\oplus\Z/2$); see \cite{cremona-1997} and the LMFDB record \cite{lmfdb-24a1}. Each of its
eight rational points pulls back to $xz=0$ or $x=z$.

\emph{$\mathcal{V}_{3/4}$.} The quartic admits a standard conic-to-Weierstrass substitution and
reduces to $E_{3/4}:y^2=x(x-300)(x-1200)$, Cremona \texttt{72a2} (conductor $72$, the same
$j$-invariant $35\,152/9$ as $E_{1/4}$, of which it is a quadratic twist, also rank $0$, torsion
$\Z/2\oplus\Z/2$); see \cite{cremona-1997,lmfdb-72a2}. Each torsion point again pulls back to a
trivial quartic solution.
\end{proof}

Lemma~\ref{lem:cremona-empty} says nothing about the diagonal $x=z$, which is where the interior
analysis has to be completed separately.

\begin{lemma}[Interior pairs with equal flanking legs]\label{lem:diagonal}
Suppose $a_i=a_{i+2}$. Then the interior pair $(i,i+1)$ carries an admissible value of
$\cos^2\theta$ if and only if $a_i=a_{i+1}=a_{i+2}$, in which case $\cos^2\theta=\tfrac14$ and
$\theta=\pi/3$.
\end{lemma}

\begin{proof}
With $x=z$ the cosine is $\cos\theta=x^2/(x^2+y^2)$, a positive rational. The values $0$ and $1$
of $\cos^2\theta$ are excluded by positivity. The value $\tfrac12$ requires
$\cos\theta=1/\sqrt2$ and the value $\tfrac34$ requires $\cos\theta=\sqrt3/2$, neither rational.
The value $\tfrac14$ requires $\cos\theta=\tfrac12$, that is $2x^2=x^2+y^2$, hence $y=x$; and
conversely $x=y=z$ gives $\cos\theta=x^2/(2x^2)=\tfrac12$.
\end{proof}

\subsection{The classification}

\begin{theorem}[Rank-two classification]\label{thm:rank2}
Let $n\ge2$ and $a\in\Q_{>0}^n$. Then $\ker\rho_a$ meets some rank-two standard parabolic
$\langle R_i,R_{i+1}\rangle$ nontrivially if and only if $e(a)+t(a)\ge1$. More precisely:
\begin{enumerate}[label=(\roman*),leftmargin=2.2em]
\item if $t(a)\ge1$, the interior pair at a triple has $\theta=\pi/3$ and
$(R_iR_{i+1})^3\in\ker\rho_a$ is an element of length $6$;
\item if $t(a)=0$ and $e(a)\ge1$, the corresponding boundary pair has $\theta=\pi/4$ and
$(R_iR_{i+1})^4\in\ker\rho_a$ is an element of length $8$, and no shorter element of any
rank-two standard parabolic lies in $\ker\rho_a$;
\item if $e(a)=t(a)=0$, no rank-two standard parabolic meets $\ker\rho_a$ nontrivially.
\end{enumerate}
\end{theorem}

\begin{proof}
Since $\langle R_i,R_{i+1}\rangle$ is an infinite dihedral standard parabolic of $W_n$, a word
$(R_iR_{i+1})^m$ is geodesic of length $2m$, and $\ker\rho_a$ meets that parabolic nontrivially
exactly when $\theta_{i,i+1}\in\pi\Q$, in which case the shortest element is
$(R_iR_{i+1})^{m}$ with $\theta_{i,i+1}=\pi\,r/m$ in lowest terms. Non-consecutive pairs
generate finite parabolics of $W_n$ itself and contribute nothing.

Suppose $\theta_{i,i+1}\in\pi\Q$ for some $i$. By the extended Niven theorem
$\cos^2\theta_{i,i+1}\in\{0,\tfrac14,\tfrac12,\tfrac34,1\}$. If the pair is a boundary pair,
Lemma~\ref{lem:boundary} gives $a_i=a_{i+1}$ with the pair at an endpoint, that is $e(a)\ge1$,
and $\theta=\pi/4$. If the pair is interior with flanking legs $x=a_i\ne a_{i+2}=z$, then
\eqref{eq:uquad} has the rational root $u=a_{i+1}^2$, so its discriminant is a rational square
and $(x,z)$ is a nontrivial point of one of $\mathcal{V}_{1/4},\mathcal{V}_{1/2},
\mathcal{V}_{3/4}$, contradicting Lemma~\ref{lem:cremona-empty}. If the pair is interior with
$a_i=a_{i+2}$, Lemma~\ref{lem:diagonal} gives $a_i=a_{i+1}=a_{i+2}$, that is $t(a)\ge1$, and
$\theta=\pi/3$. This proves (iii) and the ``only if'' direction.

Conversely, if $t(a)\ge1$ then some interior pair has $\theta=\pi/3$ by
Lemma~\ref{lem:diagonal}, the dihedral image has order $6$, and the shortest relator is
$(R_iR_{i+1})^3$ of length $6$: (i). If $t(a)=0$ and $e(a)\ge1$ then some boundary pair has
$\theta=\pi/4$ by Lemma~\ref{lem:boundary}, the dihedral image has order $8$, and the shortest
relator in that parabolic is $(R_iR_{i+1})^4$ of length $8$; by the previous paragraph no other
consecutive pair contributes an element of $\ker\rho_a$, since a boundary pair would force
$e\ge1$ with $\theta=\pi/4$ again and an interior pair would force $t\ge1$: (ii).
\end{proof}

Following the terminology of the three-dimensional analysis, call a leg tuple \emph{Class C}
when its coordinates are pairwise distinct. Class C tuples have $e=t=0$, so
Theorem~\ref{thm:rank2}(iii) applies, but so do tuples with an interior equal pair such as
$(1,2,2,3)$ and tuples with a non-adjacent repetition such as $(1,2,1,3)$; the latter is
precisely the case that Lemma~\ref{lem:diagonal}, and not Lemma~\ref{lem:cremona-empty}, has to
handle.

\begin{conjecture}[Class C faithfulness]\label{conj:master-C}
For every $n\ge3$ and every $a\in\Q_{>0}^n$ with pairwise distinct coordinates, $\rho_a$ is
injective; equivalently $u_d(a)=[t^d]\,W_n(t)$ for every $d\ge0$.
\end{conjecture}

Theorem~\ref{thm:rank2}(iii) is the rank-two half of Conjecture~\ref{conj:master-C} and is
proved. What is missing is that a nontrivial kernel element must lie in a rank-two parabolic,
which is false in general: the planar case $n=2$ is a counterexample. There the two consecutive
pairs are both boundary pairs, so for $a_1\ne a_2$ Theorem~\ref{thm:rank2}(iii) applies and no
rank-two parabolic meets $\ker\rho_a$; yet $\ker\rho_a\ne\{1\}$ by
Proposition~\ref{prop:plane}, and its shortest elements have length $20$, arising as
$w_1w_2^{-1}$ from eight pairs of distinct geodesic words $w_1,w_2$ of length $10$ using all
three generators \cite{Bonfioli1}. Conjecture~\ref{conj:master-C} is verified computationally at
$n=3$ with legs $(1,2,3)$ through depth $10$ and at $n=4$ with legs $(1,2,3,5)$ through depth
$8$ (\S\ref{sec:cert}).

\subsection{A determinant identity, and what it does not give}

The Lorentzian Gram matrix of the orthoscheme in staircase coordinates, with the standard
$(n,1)$ endpoint correction, is the tridiagonal $(n+1)\times(n+1)$ matrix $Q_n$ with
\begin{align*}
Q_n[0,0]&=1-a_1^2, & Q_n[i,i]&=a_i^2+a_{i+1}^2\ (1\le i\le n-1), & Q_n[n,n]&=1,\\
Q_n[0,1]&=a_2,     & Q_n[i,i+1]&=a_i\,a_{i+2}\ (1\le i\le n-2),   & Q_n[n-1,n]&=a_{n-1}.
\end{align*}

\begin{lemma}[Uniform Schur-complement determinant identity]\label{lem:uniform-detQ}
For every $n\ge3$ and every positive tuple $(a_1,\dots,a_n)$,
\[
\det Q_n=-\prod_{i=1}^{n}a_i^2 .
\]
\end{lemma}

\begin{proof}
Let $D_k=\det Q_n[:k{+}1,:k{+}1]$ be the leading principal minor of order $k+1$. Tridiagonal
expansion along the last row gives
\begin{equation*}\tag{$\dagger$}
D_k=Q_n[k,k]\,D_{k-1}-Q_n[k-1,k]^2\,D_{k-2}\qquad(1\le k\le n),
\end{equation*}
with $D_{-1}=1$. We claim
\begin{equation*}\tag{$\star$}
D_k=\Bigl(\prod_{i=1}^ka_i^2\Bigr)\Bigl(1-\sum_{i=1}^{k+1}a_i^2\Bigr)
\qquad\text{for }1\le k\le n-1 .
\end{equation*}

\emph{Base.} $D_0=1-a_1^2$ and $D_1=(1-a_1^2)(a_1^2+a_2^2)-a_2^2=a_1^2(1-a_1^2-a_2^2)$.

\emph{Interior step.} For $2\le k\le n-1$ the recurrence ($\dagger$) reads
$D_k=(a_k^2+a_{k+1}^2)D_{k-1}-(a_{k-1}a_{k+1})^2D_{k-2}$. Introducing
$S_{k-1}=\sum_{i=1}^{k-1}a_i^2$ and $P_{k-2}=\prod_{i=1}^{k-2}a_i^2$ and substituting ($\star$)
reduces the step to the polynomial identity
\begin{multline*}
P_{k-2}\,a_{k-1}^2\,a_k^2\,(1-S_{k-1}-a_k^2-a_{k+1}^2)
=(a_k^2+a_{k+1}^2)\,P_{k-2}\,a_{k-1}^2\,(1-S_{k-1}-a_k^2)\\
-(a_{k-1}a_{k+1})^2\,P_{k-2}\,(1-S_{k-1})
\end{multline*}
in the five independent symbols $(a_{k-1},a_k,a_{k+1},S_{k-1},P_{k-2})$. It is independent of
$k$ and of $n$.

\emph{Final step.} At $k=n$ the recurrence uses $Q_n[n,n]=1$ and $Q_n[n-1,n]^2=a_{n-1}^2$,
giving $D_n=D_{n-1}-a_{n-1}^2D_{n-2}$. With $S_{n-1}=\sum_{i=1}^{n-1}a_i^2$ and
$P_{n-2}=\prod_{i=1}^{n-2}a_i^2$,
\[
D_{n-1}-a_{n-1}^2D_{n-2}
=P_{n-2}\,a_{n-1}^2\bigl[(1-S_{n-1}-a_n^2)-(1-S_{n-1})\bigr]
=-P_{n-2}\,a_{n-1}^2a_n^2=-\prod_{i=1}^na_i^2 .
\]
The three identities close the induction for every $n\ge3$, with no per-$n$ verification. All of
them, the induction that assembles them, and the boundary step are machine-checked in Lean~4
over an arbitrary commutative ring; no positivity, ordering or characteristic hypothesis is used
(\S\ref{sec:cert}).
\end{proof}

\begin{remark}[Why nonvanishing of $\det Q_n$ does not give faithfulness]\label{rem:detQ-nonsequitur}
An earlier version of this material argued from Lemma~\ref{lem:uniform-detQ} that
$\det Q_n\ne0$, hence the form is nondegenerate, hence $\rho_a$ is faithful by Tits' theorem.
That argument is a non sequitur, in three independent ways, and the deduction is withdrawn.

First, it proves too much. The hypothesis of Lemma~\ref{lem:uniform-detQ} is that the legs are
positive and nothing more, so the conclusion $\det Q_n\ne0$ holds at every positive tuple,
including $(1,1,\dots,1)$. Were the inference valid, $\rho_a$ would be faithful there too. It is
not: breadth-first search at $n=3$ with legs $(1,1,1)$ gives layer counts $1,4,9,17,28,42$
against the envelope $1,4,9,18,36,72$, a deficit already at depth $3$, in agreement with
Theorem~\ref{thm:cd-general}, which assigns $\cd_3=3$ to that tuple. The same computation at
$n=4$ with legs $(1,1,1,1)$ gives $1,5,14,31,59,101$ against $1,5,14,33,75,169$.

Second, Tits' theorem concerns the canonical geometric representation of an abstract Coxeter
system, on the space spanned by the simple roots with the form
$B(\alpha_i,\alpha_j)=-\cos(\pi/m_{ij})$ \cite{Bourbaki}. It is not a statement about $\rho_a$,
which is a different representation, on $\R^n$ rather than $\R^{n+1}$ and affine rather than
linear. Nor is $Q_n$ a Coxeter Gram matrix: its diagonal is not identically $1$, and
$Q_n[0,0]=1-a_1^2$ is negative whenever $a_1>1$.

Third, run at $n=2$ the argument returns a false answer. There $\det Q_2\ne0$ just as well, so it
would give $u_d(a)=[t^d]W_2(t)=F_{d+3}$ for all $d$; in fact the two agree only up to $d=9$ and
differ at $d=10$, where the orbit count is $225$ against the envelope's $233$, a deficit of $8$
accounted for by eight coincidences between distinct words of length $10$ \cite{Bonfioli1}.
\end{remark}

\section{The collision depth}\label{sec:cd}

\begin{definition}\label{def:cd}
For $a\in\Q_{>0}^n$ the \emph{collision depth} is
\[
\cd_n(a)=\inf\bigl\{d\ge0:u_d(a)<[t^d]\,W_n(t)\bigr\}\in\{1,2,\dots\}\cup\{\infty\}.
\]
\end{definition}

By Corollary~\ref{cor:cd-is-K}, $\cd_n(a)=K(a)/2$, and by Lemma~\ref{lem:no-length-four},
$K(a)\ge6$. Conjecture~\ref{conj:master-C} says exactly that $\cd_n=\infty$ for every Class C
tuple. The next theorem describes $\cd_n$ for every leg multiplicity pattern. Although the
multiplicity partition of $(a_1,\dots,a_n)$ has $p(n)$ possible shapes, the collision depth
takes only three values, determined by the two local positional statistics $(e,t)$.

\begin{theorem}[Collision depth, all $n\ge3$]\label{thm:cd-general}
Let $n\ge3$ and $a\in\Q_{>0}^n$. Then $\cd_n(a)\ge3$, and
\[
\cd_n(a)=
\begin{cases}
\infty & \text{if }e=t=0\ \text{ and }\rho_a\text{ is injective},\\
4      & \text{if }e\ge1,\ t=0\ \text{ and }K(a)\ne6,\\
3      & \text{if }t\ge1 .
\end{cases}
\]
The third case is unconditional. In the second case $\cd_n(a)\in\{3,4\}$ unconditionally, and
$K(a)=6$ would require an element of $\ker\rho_a$ of length $6$ lying in no rank-two standard
parabolic; no such element occurs in any tuple of Table~\ref{tab:atlas}. In the first case
$e=t=0$ includes every Class C tuple, so that branch is $\infty$ exactly when
Conjecture~\ref{conj:master-C} holds.
\end{theorem}

\begin{proof}
$\cd_n(a)=K(a)/2$ with $K(a)$ even and at least $6$, by Corollary~\ref{cor:cd-is-K} and
Lemma~\ref{lem:no-length-four}; this gives $\cd_n(a)\ge3$ in all cases.

If $t\ge1$, Theorem~\ref{thm:rank2}(i) exhibits an element of $\ker\rho_a$ of length $6$, so
$K(a)\le6$, hence $K(a)=6$ and $\cd_n(a)=3$.

If $e\ge1$ and $t=0$, Theorem~\ref{thm:rank2}(ii) exhibits an element of length $8$, so
$6\le K(a)\le8$ and $K(a)$ is even, giving $\cd_n(a)\in\{3,4\}$. By Theorem~\ref{thm:rank2}(ii)
no element of a rank-two standard parabolic lies in $\ker\rho_a$ with length below $8$, so
$K(a)=6$ is possible only through an element of length $6$ lying in no rank-two standard
parabolic. Absent such an element, $K(a)=8$ and $\cd_n(a)=4$.

If $e=t=0$ then $\cd_n(a)=\infty$ if and only if $\ker\rho_a=\{1\}$, by
Corollary~\ref{cor:alternative}.
\end{proof}

\begin{remark}[What the second case leaves open]\label{rem:second-case}
The gap in the second case is genuine and is a special case of the general problem: the rank-two
analysis of \S\ref{sec:rank2} controls elements of $\ker\rho_a$ that lie in a dihedral standard
parabolic, and nothing in this paper controls the others. The gap is narrow: only the single
value $K(a)=6$ is at stake, and only for tuples with an endpoint equal pair and no triple.
\end{remark}

\begin{remark}[The run-length rule fails in $n\ge4$]\label{rem:R-rule}
A natural but incorrect heuristic is to declare $\cd_n$ finite as soon as some consecutive equal
pair appears, that is as soon as the run length $R=\max_i\#\{j:a_j=a_i\}$ is at least $2$. At
$n=3$ this rule coincides with the $(e,t)$ rule, because every consecutive equal pair in a
three-term sequence is automatically at an endpoint. At $n\ge4$ the two diverge: the orthoscheme
$(1,2,2,3)$ has an interior equal pair, so $R=2$ but $(e,t)=(0,0)$, and
Theorem~\ref{thm:cd-general} places it in the first case; empirically its layer counts
$1,5,14,33,75,169,\dots$ are identical to those of the Class C tuple $(1,2,3,5)$ through every
computed depth, falsifying the run-length rule.
\end{remark}

\begin{table}[ht!]
\centering
\small
\renewcommand{\arraystretch}{1.2}
\begin{tabular}{c l c c l}
\toprule
$n$ & leg sequence & $(e,t)$ & $\cd_n$ & first six layer counts \\
\midrule
3 & $(2,3,5)$          & $(0,0)$ & $\infty$ & $1,4,9,18,36,72$ \\
3 & $(1,1,2)$          & $(1,0)$ & $4$      & $1,4,9,18,35,67$ \\
3 & $(1,1,1)$          & $(2,1)$ & $3$      & $1,4,9,17,28,42$ \\
\midrule
4 & $(1,2,3,5)$        & $(0,0)$ & $\infty$ & $1,5,14,33,75,169$ \\
4 & $(1,2,2,3)$        & $(0,0)$ & $\infty$ & $1,5,14,33,75,169$ \\
4 & $(1,1,2,3)$        & $(1,0)$ & $4$      & $1,5,14,33,74,163$ \\
4 & $(1,1,2,2)$        & $(2,0)$ & $4$      & $1,5,14,33,73,157$ \\
4 & $(1,1,1,2)$        & $(1,1)$ & $3$      & $1,5,14,32,67,135$ \\
\midrule
5 & $(1,2,3,5,7)$      & $(0,0)$ & $\infty$ & $1,6,20,54,136,334$ \\
5 & $(1,2,2,3,5)$      & $(0,0)$ & $\infty$ & $1,6,20,54,136,334$ \\
5 & $(1,1,2,3,5)$      & $(1,0)$ & $4$      & $1,6,20,54,135,327$ \\
5 & $(2,1,1,1,2)$      & $(0,1)$ & $3$      & $1,6,20,53,128,297$ \\
\midrule
6 & $(1,1,2,3,2,2)$    & $(2,0)$ & $4$      & $1,7,27,82,224,581$ \\
\bottomrule
\end{tabular}
\caption{A small atlas of representative orthoschemes in $n=3,4,5,6$ illustrating each
attainable $(e,t)$ class. The layer counts are exact-rational breadth-first search
(\S\ref{sec:cert}); in every row the observed collision depth is the value predicted by
Theorem~\ref{thm:cd-general}, and the entries $\infty$ are observed agreement with the envelope
as far as computed, not a proof. The two $n=4$ rows with $(e,t)=(0,0)$ have identical counts
through every computed depth despite differing in run structure.}
\label{tab:atlas}
\end{table}

The rows with $\cd_n=\infty$ saturate the upper bound $[t^d]\,W_n(t)$ as far as they are
computed, as Conjecture~\ref{conj:master-C} predicts; for $n=3$ this reproduces
$u_d=9\cdot2^{d-2}$ for $d\ge2$.

\section{The plane, and why its argument does not extend}\label{sec:plane}

\begin{proposition}[The plane]\label{prop:plane}
For $n=2$ and every $a\in U$, $\rho_a$ is not injective, so alternative (ii) of
Corollary~\ref{cor:alternative} holds.
\end{proposition}

\begin{proof}
$W(a)\le\Isom(\R^2)=\R^2\rtimes O(2)$. The group $O(2)$ contains the abelian $SO(2)$ with index
$2$, so $\Isom(\R^2)$ is virtually solvable as an abstract group, hence amenable, and so is
every subgroup of it; in particular $W(a)$ is amenable. The right-angled Coxeter group $W_2$ is
not amenable: its defining graph $\Gamma_2$ is not complete multipartite, since its complement
$P_3$ is connected and is not a disjoint union of cliques, and a right-angled Coxeter group
whose defining graph is not complete multipartite contains a free subgroup of rank $2$
\cite{Davis}. An amenable group is not isomorphic to a non-amenable one, so $\rho_a$ is not
injective.
\end{proof}

\begin{corollary}[The planar horizon]\label{cor:plane}
For $n=2$ and every shape admissible in the sense of \cite{Bonfioli1}, the collision depth is
$\cd_2=10$, equivalently $K(a)=20$: $u_d(a)=F_{d+3}=[t^d]W_2(t)$ for $1\le d\le9$ (while $u_0=1$
and $F_3=2$), and at $d=10$ the count is $225$ against the envelope's $233$. Past that horizon
the sequence is \textup{A396406}, with growth rate $\beta_2=1.4916\ldots<\varphi=r_2$, and its
generating function is transcendental over $\Q(x)$ \cite{Bonfioli1}. This is the only dimension
in which the deviation has been computed past its onset.
\end{corollary}

The isoceles right triangle $a_1=a_2$ is excluded from that statement, and must be: there
$e(a)=1$, so Theorem~\ref{thm:rank2}(ii) gives an element of $\ker\rho_a$ of length $8$ and
hence $\cd_2\le4$.

\begin{remark}[The amenability trap: why the planar argument stops at $n=2$]
\label{rem:amenability-fails}
The proof of Proposition~\ref{prop:plane} does not extend to $n\ge3$, and the natural extension
is false. One is tempted to argue that $\Isom(\R^n)=\R^n\rtimes O(n)$ is amenable because $O(n)$
is compact. Compactness gives amenability of $O(n)$ as a \emph{topological} group, with respect
to Haar measure, and that property does not pass to abstract subgroups. As an abstract discrete
group $O(n)$ is non-amenable for every $n\ge3$: already $SO(3)$ contains a free subgroup of rank
$2$ (Hausdorff; see \cite[Thm.~2.1]{Wagon}), which is the standard ingredient of the
Banach--Tarski paradox. Hence $\Isom(\R^n)$ is non-amenable as an abstract group for $n\ge3$ and
yields no obstruction at all. At $n=2$ the argument survives only because $O(2)$ is virtually
abelian.
\end{remark}

\section{What is proved, and the open problems}\label{sec:open}

Table~\ref{tab:status} lists the status of every numbered statement. Four things are proved
unconditionally and dimension-freely: the dichotomy (Theorem~\ref{thm:dichotomy}); generic
rigidity (Theorem~\ref{thm:rigidity}), which makes ``the generic $n$-orthoscheme growth'' a
well-defined object realised by a co-null set of shapes; rationality of the envelope with rate
$r_n$ (Theorem~\ref{thm:envelope}); and the rank-two classification
(Theorem~\ref{thm:rank2}), with its consequences that $\cd_n\ge3$ for every leg tuple and that
$\cd_n=3$ at every tuple with three consecutive equal legs. Whether $\cd_n=3$ can also occur
without such a triple is not decided here; it would require a kernel element of length $6$ in no
rank-two parabolic (Remark~\ref{rem:second-case}).

Which alternative of Corollary~\ref{cor:alternative} holds is settled only in the plane, by
Proposition~\ref{prop:plane}, where it is (ii) with horizon $10$. For $3\le n\le7$ the geometric
growth has been computed to depth $13$ to $16$ and agrees with the envelope throughout that
range. That agreement is consistent with both alternatives: under (i) it would continue forever,
and under (ii) it would end at a horizon beyond the computed range. Nothing here decides between
them.

\begin{problem}[The alternative in dimension $n\ge3$]\label{prob:alternative}
Fix $n\ge3$.
\begin{enumerate}[label=(\alph*),leftmargin=2.2em]
\item \emph{Generic form.} Is $N_{\mathrm{gen}}=\bigcap_{a\in U}\ker\rho_a$ trivial?
Equivalently, by Theorem~\ref{thm:rigidity}, is $\rho_a$ injective for every $a$ in the co-null
set $U_{\mathrm{gen}}$; equivalently, by Corollary~\ref{cor:alternative}, is the growth series
of the generic $n$-orthoscheme group the rational envelope $W_n(t)$ in every degree?
\item \emph{Arithmetic form.} Is $\rho_a$ injective for every Class C tuple
$a\in\Q_{>0}^n$? This is Conjecture~\ref{conj:master-C}.
\end{enumerate}
\end{problem}

The two forms are related in one direction only. By Corollary~\ref{cor:one-shape}, a single
faithful tuple, rational or not, answers (a) affirmatively; so (b) implies (a). The converse
fails to be automatic because $U_{\mathrm{gen}}$ contains no tuple with algebraically dependent
coordinates, hence no rational tuple: an affirmative answer to (a) says nothing directly about
any given rational $a$. If the answer to (a) is affirmative for all $n\ge3$, the plane is the
unique dimension with non-rational orbit growth. If instead $\ker\rho_a$ is nontrivial in some
dimension $n\ge3$, the growth deviates below the envelope at $\cd_n(a)=K(a)/2$, and the
arithmetic of the deviated series past that horizon is open even in the cases where the horizon
can be located.

Three obstructions are worth stating precisely, because they close off the routes that look
available.

\begin{enumerate}[label=(O\arabic*),leftmargin=2.6em,itemsep=0.4em]
\item \emph{The rank-two half is all that the Diophantine method gives.} Faithfulness would
follow from Theorem~\ref{thm:rank2}(iii) if every nontrivial kernel element lay in a rank-two
parabolic. That implication is false in general: at $n=2$ with $a_1\ne a_2$ no rank-two parabolic
meets $\ker\rho_a$, yet $\ker\rho_a$ is nontrivial, its shortest elements having length $20$ and
using all three generators \cite{Bonfioli1}. So no soft upgrade of Theorem~\ref{thm:rank2} to
Conjecture~\ref{conj:master-C} is available, and the same obstruction is what leaves the value
$K(a)=6$ unexcluded in the second case of Theorem~\ref{thm:cd-general}
(Remark~\ref{rem:second-case}).

\item \emph{The amenability trap.} The one soft argument that settles $n=2$ is unavailable for
$n\ge3$, and not merely unproved: $O(n)$ contains a free subgroup of rank $2$ as an abstract
group for every $n\ge3$, so $\Isom(\R^n)$ is not amenable as an abstract group and imposes no
constraint whatsoever on its subgroups' isomorphism type
(Remark~\ref{rem:amenability-fails}).

\item \emph{Finite presentation.} At $n=2$, Isola and Piergallini~\cite{isola-piergallini} show
that the generic triangle reflection group is not finitely presented, which gives non-injectivity
independently of (O2). We do not know whether the generic $n$-orthoscheme reflection group is
finitely presented for $n\ge3$; an answer either way would decide
Problem~\ref{prob:alternative}(a), since $W_n$ is finitely presented.
\end{enumerate}

Three further questions are open and are stated for the record.

\begin{enumerate}[label=(Q\arabic*),leftmargin=2.6em,itemsep=0.35em]
\item \emph{Other diagrams.} Does an analogue of Theorem~\ref{thm:cd-general} hold for other
families of simplicial reflection chambers in $\R^n$, for instance Coxeter orthoschemes with
bent paths or simplices from non-path-shaped diagrams? The analysis here uses the
right-angled-with-one-nontrivial-bond structure of the path diagram throughout.

\item \emph{Intermediate rates.} In every dimension $n\ge3$, multiplicity patterns that are
neither Class C nor saturated by triples produce growth rates strictly between two consecutive
values of $r_n$. Whether these intermediate rates are algebraic over $\Q$ at all is open.

\item \emph{Beyond $(e,t)$.} Theorem~\ref{thm:cd-general} shows that the collision depth is
determined by the positional statistics $(e,t)$. Is there a finer positional invariant that
determines the entire sequence $u_d(a)$, not only its first deviation?
\end{enumerate}

\begin{table}[ht!]
\centering
\footnotesize
\renewcommand{\arraystretch}{1.15}
\begin{tabular}{l l l}
\toprule
Statement & Status & Note \\
\midrule
Lemma~\ref{lem:normals} (facet normals)          & proved (Lean) & \texttt{OrthoschemeNormals} \\
Theorem~\ref{thm:rigidity} (generic rigidity)    & proved                  & \\
Corollary~\ref{cor:one-shape}                    & proved                  & \\
Theorem~\ref{thm:dichotomy} (dichotomy)          & proved (Lean) & \texttt{EnvelopeDichotomy} \\
Corollaries~\ref{cor:alternative}, \ref{cor:cd-is-K} & proved              & \\
Lemma~\ref{lem:no-length-four} ($K\ge6$)         & proved                  & \\
Theorem~\ref{thm:envelope} (envelope)            & proved (Lean, partly) & \texttt{EnvelopeDichotomy}, \texttt{NdimRate} \\
Corollary~\ref{cor:asymptotic}                   & proved (Lean) & \texttt{NdimRate} \\
Lemma~\ref{lem:boundary} (boundary pairs)        & proved (Lean) & \texttt{RankTwoExclusion} \\
Lemma~\ref{lem:cremona-empty} (three quartics)   & proved                  & rank-$0$ descents, not formalised \\
Lemma~\ref{lem:diagonal} (equal flanking legs)   & proved                  & \\
Theorem~\ref{thm:rank2} (rank-two classification)& proved                  & interior half uses Lemma~\ref{lem:cremona-empty} \\
Lemma~\ref{lem:uniform-detQ} ($\det Q_n$)        & proved (Lean) & \texttt{OrthoschemeDet} \\
Theorem~\ref{thm:cd-general}, case $t\ge1$       & proved                  & \\
Theorem~\ref{thm:cd-general}, case $e\ge1,t=0$   & $\cd_n\in\{3,4\}$ proved & value $4$ conditional, Rem.~\ref{rem:second-case} \\
Theorem~\ref{thm:cd-general}, case $e=t=0$       & conditional             & equivalent to Conj.~\ref{conj:master-C} \\
Conjecture~\ref{conj:master-C} (Class C)         & conjecture              & checked to depth $10$ ($n=3$), $8$ ($n=4$) \\
Proposition~\ref{prop:plane} (the plane)         & proved                  & \\
Problem~\ref{prob:alternative}                   & open in both directions & \\
\bottomrule
\end{tabular}
\caption{Status of the numbered statements. ``Lean'' means that the indicated Lean~4 file, which
builds against Mathlib with no \texttt{sorry}, covers the statement; ``partly'' means that it
covers some steps of the proof and not others. The scope of each file is described in
\S\ref{sec:cert}.}
\label{tab:status}
\end{table}

\section{Related work}\label{sec:related}

\paragraph{Path Coxeter groups and growth-rate calculus.}
The closed form $r_n=1+2\cos(2\pi/(n+3))$ is the signature of a linear, path-shaped Coxeter
group whose Poincar\'e-series denominator reduces to a tridiagonal Chebyshev-type
characteristic polynomial. Cannon and Wagreich~\cite{cannon-wagreich}, Floyd and
Plotnick~\cite{floyd-plotnick} and Parry~\cite{parry-1989} developed the growth-rate calculus of
such groups, and Kellerhals and Perren~\cite{kellerhals-perren} treat right-angled Coxeter growth
systematically. Theorem~\ref{thm:envelope} fits into that framework with parameter $m=n+3$.

\paragraph{A Salem-type family on the surface-dynamics side.}
The family $\{1+2\cos(2\pi/m)\}_{m\ge3}$ to which $r_n$ belongs also arises in birational
surface dynamics. Diller and Favre~\cite{diller-favre} introduced the dynamical-degree framework
for bimeromorphic maps of surfaces; McMullen~\cite{mcmullen-2007} showed that this family arises
as dynamical degrees of automorphisms of blowups of $\mathbb{P}^2$ associated with the $E_n$
T-shaped Coxeter graph, which is not the path graph occurring here; and
Bedford and Kim~\cite{bedford-kim} realised specific members through linear-fractional
recurrences. The rate $r_n$ lies in this family, but the right-corner orthoscheme reflection
group is not known to be realised as a McMullen involution product on a rational surface, and
the coincidence is at the level of the algebraic-number family only. A birational-geometry
construction realising $r_n$ as a dynamical degree on a $\mathbb{P}^2$-blowup would be of
independent interest.

\section{Certification}\label{sec:cert}

\paragraph{Orbit counts.}
The layer counts quoted above, the envelope coefficients and the Diophantine searches are
produced by \texttt{code/zeta\_probe/tools/ortho\_cd} (Rust, exact arithmetic in $\mathbb{F}_p$
at two primes). It builds the $n+1$ facet reflections directly from the legs, enumerates the
group by breadth-first search, and compares layer by layer against $(1+t)^{n+1}/J_{n+1}(t)$. In
particular it produces the counts of Remark~\ref{rem:detQ-nonsequitur} at $(1,1,1)$ and
$(1,1,1,1)$ and at $n=2$, confirms the $(e,t)$ prediction on every tuple of
Table~\ref{tab:atlas} including the non-adjacent repetition $(1,2,1,3)$, and searches the three
quartics of Lemma~\ref{lem:cremona-empty} for nontrivial solutions with legs up to $300$,
finding none, while exhibiting rational points of the corresponding conics in $(X,Z)$ to show
that the restriction to squares is not cosmetic.

\paragraph{The envelope.}
Theorem~\ref{thm:envelope} is certified by \texttt{code/zeta\_probe/tools/racg\_envelope}
(Rust, exact integer polynomial arithmetic), which for $2\le n\le30$ checks that the recurrence
for $J_m$ and the Steinberg expansion agree coefficient by coefficient, performs the division by
$(1+t)^{\lfloor(n+1)/2\rfloor}$ exactly, and matches the smallest positive root of $K_{n+1}$
against $1/\bigl(1+2\cos\frac{2\pi}{n+3}\bigr)$. Independently, for $2\le n\le6$ it enumerates
$W_n$ by breadth-first search in the canonical geometric representation, whose matrices are
integral here because the bilinear form takes only the values $1,0,-1$ and which is faithful by
Tits' theorem \cite{Bourbaki}; the sphere sizes agree with the series coefficients to depths $9$
to $14$. The same program records two checks on the shape of the answer: the variant substitution
$-1/(1+t)$ in place of $-t/(1+t)$ gives a different rational function already at $n=2$, and the
polynomial $1-2t-3t^2+4t^3-t^4$ is not the $n=7$ denominator, being the negated reciprocal of
$1-4t+3t^2+2t^3-t^4$ and predicting $6$ elements at depth $1$ in a group with $8$ generators.

\paragraph{Lean 4.}
Five files in \texttt{lean/with\_mathlib/} build against Mathlib with no \texttt{sorry}.

\begin{center}\footnotesize
\renewcommand{\arraystretch}{1.15}
\begin{tabular}{ll}
\toprule
Paper result & Lean declaration \\
\midrule
\multicolumn{2}{l}{\emph{\texttt{OrthoschemeNormals.lean}} (7 theorems)}\\
Lemma~\ref{lem:normals}, support pattern & \texttt{supp\_disjoint}\\
Lemma~\ref{lem:normals}, orthogonality on disjoint supports
  & \texttt{mul\_eq\_zero\_of\_disjoint}, \texttt{sum\_eq\_zero\_of\_disjoint}\\
Lemma~\ref{lem:normals}, the three adjacent products
  & \texttt{dot\_zero\_one}, \texttt{dot\_interior}, \texttt{dot\_last}\\
Lemma~\ref{lem:normals}, combined & \texttt{orthogonality\_pattern}\\
\addlinespace
\multicolumn{2}{l}{\emph{\texttt{EnvelopeDichotomy.lean}} (12 theorems)}\\
Thm.~\ref{thm:dichotomy}, horizon arithmetic
  & \texttt{two\_mul\_lt\_of\_lt\_horizon}, \texttt{split\_lengths}\\
Thm.~\ref{thm:dichotomy}, even case & \texttt{card\_image\_lt\_of\_collision}\\
Thm.~\ref{thm:dichotomy}, odd case & \texttt{card\_lt\_of\_escapes}\\
Thm.~\ref{thm:envelope}, three-term recurrence
  & \texttt{cos\_three\_term}, \texttt{sin\_three\_term}, \texttt{J\_recurrence}\\
Thm.~\ref{thm:envelope}, root condition
  & \texttt{cot\_add\_cot\_eq\_zero\_iff}, \texttt{J\_eq\_zero\_iff}\\
Thm.~\ref{thm:envelope}, recovering $t$ from $\theta$
  & \texttt{t\_of\_theta}, \texttt{rate\_of\_theta}\\
Remark~\ref{rem:notpath} & \texttt{rate\_ne\_path\_spectral\_radius}\\
\addlinespace
\multicolumn{2}{l}{\emph{\texttt{NdimRate.lean}} (7 theorems)}\\
Remark~\ref{rem:pell}, Pascal step & \texttt{pascal\_step} (no axioms), \texttt{coeff\_pascal}\\
Remark~\ref{rem:pell}, Binet formula & \texttt{sq\_eq}, \texttt{pell\_binet}\\
Remark~\ref{rem:pell}, root translation & \texttt{root\_translation}\\
Corollary~\ref{cor:asymptotic} & \texttt{rate\_residual}, \texttt{three\_sub\_r}\\
\addlinespace
\multicolumn{2}{l}{\emph{\texttt{RankTwoExclusion.lean}} (10 theorems)}\\
Lemma~\ref{lem:boundary}, the five Niven values
  & \texttt{niven\_zero}, \texttt{niven\_one}, \texttt{niven\_quarter},\\
  & \texttt{niven\_three\_quarter}, \texttt{niven\_half},
    \texttt{boundary\_no\_relation}\\
Lemma~\ref{lem:boundary}, $3$ is not a rational square
  & \texttt{sq\_ne\_three}, \texttt{not\_isSquare\_three\_int}\\
Boundary cosine nonzero, and $\cos^2\ne\tfrac14$ there
  & \texttt{no\_length\_two}, \texttt{no\_length\_three\_boundary}\\
\addlinespace
\multicolumn{2}{l}{\emph{\texttt{OrthoschemeDet.lean}} (3 theorems)}\\
Lemma~\ref{lem:uniform-detQ}, closed form ($\star$) & \texttt{minor\_closed\_form}\\
Lemma~\ref{lem:uniform-detQ}, boundary step & \texttt{det\_Q}\\
Lemma~\ref{lem:uniform-detQ}, full statement & \texttt{det\_Q\_of\_recurrence}\\
\bottomrule
\end{tabular}
\end{center}

\medskip
\texttt{NdimRate.lean} and \texttt{OrthoschemeDet.lean} end with \verb|#print axioms|. Of the
ten declarations they contribute to the table, six report exactly the standard axioms
\verb|propext|, \verb|Classical.choice|, \verb|Quot.sound|, and four report strictly fewer, with
\texttt{pascal\_step} depending on no axioms at all.

\medskip
\noindent\textbf{What is not formalised.} \texttt{RankTwoExclusion.lean} covers the boundary
half of Theorem~\ref{thm:rank2} only; the interior cases rest on the rank-$0$ Mordell descents
of Lemma~\ref{lem:cremona-empty} and are not formalised. \texttt{EnvelopeDichotomy.lean}
formalises the arithmetic of Theorem~\ref{thm:dichotomy} and the trigonometry of the root
identity, but no word metric, no Coxeter group and no clique polynomial: the group theory around
those steps is on paper. The step identifying the growth rate with the largest root $t_1$ is a
dominant-pole argument on a rational generating function, proved on paper; so is the induction
assembling the Pascal step into $c_{\Gamma_n}=P_{n+3}$ in Remark~\ref{rem:pell}. Everything
downstream of those two inputs is machine-checked. Lemma~\ref{lem:no-length-four},
Lemma~\ref{lem:diagonal} and Theorem~\ref{thm:rigidity} are paper proofs.

\section*{Acknowledgments}

Computational assistance from large language model assistants was used at multiple stages of
this work; all mathematical claims have been verified by the author, who is solely responsible
for the content of this paper and for any errors.

\begin{thebibliography}{99}

\bibitem{bedford-kim}
E.~Bedford, K.~Kim, \emph{Dynamics of rational surface automorphisms: linear fractional
recurrences}, J.\ Geom.\ Anal.\ \textbf{19} (2009), no.~3, 553--583.

\bibitem{Bonfioli1}
V.~Bonfioli, \emph{The universal right-triangle reflection sequence}, companion paper; code and
formalization at \texttt{github.com/ElVec1o/pythagorean-reflection-sequence}.

\bibitem{Bourbaki}
N.~Bourbaki, \emph{Groupes et alg\`ebres de Lie, Chapitres IV--VI}, Hermann, Paris, 1968.
(Geometric representation of a Coxeter group and its faithfulness; Steinberg's identity.)

\bibitem{cannon-wagreich}
J.~W. Cannon, P.~Wagreich, \emph{Growth functions of surface groups}, Math.\ Ann.\ \textbf{293}
(1992), no.~2, 239--257.

\bibitem{cassels-1991}
J.~W.~S. Cassels, \emph{Lectures on Elliptic Curves}, London Mathematical Society Student Texts
24, Cambridge University Press, 1991.

\bibitem{conway-jones}
J.~H. Conway, A.~J. Jones, \emph{Trigonometric Diophantine equations (On vanishing sums of roots
of unity)}, Acta Arith.\ \textbf{30} (1976), 229--240.

\bibitem{Coxeter1934}
H.~S.~M. Coxeter, \emph{Discrete groups generated by reflections}, Ann.\ of Math.\ (2)
\textbf{35} (1934), 588--621.

\bibitem{cremona-1997}
J.~E. Cremona, \emph{Algorithms for Modular Elliptic Curves}, 2nd edition, Cambridge University
Press, 1997.

\bibitem{Davis}
M.~W. Davis, \emph{The Geometry and Topology of Coxeter Groups}, London Math.\ Soc.\ Monographs
\textbf{32}, Princeton University Press, 2008. (Right-angled Coxeter groups; finite parabolics
as cliques; the complete multipartite criterion.)

\bibitem{diller-favre}
J.~Diller, C.~Favre, \emph{Dynamics of bimeromorphic maps of surfaces}, Amer.\ J.\ Math.\
\textbf{123} (2001), no.~6, 1135--1169.

\bibitem{floyd-plotnick}
W.~J. Floyd, S.~P. Plotnick, \emph{Growth functions on Fuchsian groups and the Euler
characteristic}, Invent.\ Math.\ \textbf{88} (1987), no.~1, 1--29.

\bibitem{humphreys90}
J.~E. Humphreys, \emph{Reflection Groups and Coxeter Groups}, Cambridge Studies in Advanced
Mathematics 29, Cambridge University Press, 1990.

\bibitem{isola-piergallini}
S.~Isola, R.~Piergallini, \emph{On the generic triangle group}, preprint, arXiv:1405.1881.

\bibitem{kellerhals-perren}
R.~Kellerhals, G.~Perren, \emph{On the growth of cocompact hyperbolic Coxeter groups},
European J.\ Combin.\ \textbf{32} (2011), 1299--1316.

\bibitem{lmfdb-24a1}
LMFDB Collaboration, elliptic curve \texttt{24a1} over $\Q$,
\url{https://www.lmfdb.org/EllipticCurve/Q/24/a/1}.

\bibitem{lmfdb-72a2}
LMFDB Collaboration, elliptic curve \texttt{72a2} over $\Q$,
\url{https://www.lmfdb.org/EllipticCurve/Q/72/a/2}.

\bibitem{mcmullen-2007}
C.~T. McMullen, \emph{Dynamics on blowups of the projective plane}, Publ.\ Math.\ Inst.\ Hautes
\'Etudes Sci.\ \textbf{105} (2007), 49--89.

\bibitem{mordell-1969}
L.~J. Mordell, \emph{Diophantine Equations}, Pure and Applied Mathematics 30, Academic Press,
1969.

\bibitem{parry-1989}
W.~Parry, \emph{Growth series of some hyperbolic Coxeter groups}, Geom.\ Dedicata \textbf{32}
(1989), 75--96.

\bibitem{Steinberg}
R.~Steinberg, \emph{Endomorphisms of Linear Algebraic Groups}, Mem.\ Amer.\ Math.\ Soc.\
\textbf{80} (1968). (The Poincar\'e-series identity for Coxeter groups.)

\bibitem{Wagon}
S.~Wagon, \emph{The Banach--Tarski Paradox}, Cambridge University Press, 1985. (Free subgroups
of $SO(3)$.)

\end{thebibliography}

\end{document}
